\chapter{Quaternions}
Quaternions are four dimensional complex numbers that is defined as the following:

\begin{center}
    $Q = w + a\mathbf{i} + b\mathbf{j} + c\mathbf{k}=\begin{bmatrix}w & \begin{pmatrix}a & b & c\end{pmatrix}\end{bmatrix}=\begin{bmatrix}w & \vec{V}\end{bmatrix}$
\end{center}

where:
\begin{itemize}
    \item $w, a, b, c$: are the real numbers $\mathbb{R}$ that denotes the weight of each axis in the complex number.
    \item $\mathbf{i}, \mathbf{j}, \mathbf{k}$: are the axes representing the 2nd up to 4th dimensions. ($w$ is it's own dimension as well! In the real dimension :3)
    \item satisfies $\mathbf{i}^2=\mathbf{j}^2=\mathbf{k}^2=\mathbf{ijk}=-1$
    \item $\vec{V}$: is referred to as the \textbf{vector component} of a quaternion.
    \item $w$: is referred to as the \textbf{scalar component} of a quaternion.
\end{itemize}



\subsection{Quaternion Properties}

\textbf{Magnitude} of a quaternion:
\begin{center}
    $\|Q\| = \|\begin{bmatrix}w & \vec{V}\end{bmatrix}\| = \sqrt{w^2 + \|\vec{V}\|^2}$
\end{center}

If we define $\vec{V} = \begin{bmatrix}a & b & c\end{bmatrix}$ then:
\begin{center}
    $\|Q\| = \sqrt{w^2 + a^2 + b^2 + c^2}$
\end{center}

Which is similar to a four dimensional vector's eucledian / pythagorean distance.

If the magnitude of a quaternion is $1$, then we call it a \textbf{unit quaternion}(similar to a unit vector).

A unit quaternion is special, as it represents "pure rotation", or "rotational displacement" for
three-dimensional spaces.

\sgap
\textbf{Identity} quaternion:
\begin{center}
    $I = \begin{bmatrix}1 & \vec{0}\end{bmatrix}$
\end{center}

The identity quaternion is special because it means "\textbf{no rotational displacement}", it
is particularly important with quaternion inverses and multiplication (see \textbf{Quaternion Operations} section for more detail).


\subsection{Quaternion Operations}
\textbf{Scalar Multiplication} of a quaternion:
\begin{center}
    $kQ = k \begin{bmatrix}w & \vec{V}\end{bmatrix} =\begin{bmatrix}kw & k\vec{V}\end{bmatrix}$
\end{center}


\textbf{Scalar Division} of a quaternion:
\begin{center}
    $\frac{Q}{k} = \frac{\begin{bmatrix}w & \vec{V}\end{bmatrix}}{k} =\begin{bmatrix}\frac{w}{k} & \frac{1}{k}\vec{V}\end{bmatrix}$
\end{center}

\sgap
\textbf{Multiplication} between two quaternions:
\begin{center}
    \begin{flalign}
        Q_1Q_2 & = \begin{bmatrix}w_1 & \vec{V_1}\end{bmatrix}\begin{bmatrix}w_2 & \vec{V_2}\end{bmatrix} \\
            & = \begin{bmatrix}w_1w_2 - \vec{V_1} \cdot \vec{V_2} & w_1\vec{V_2} + w_2\vec{V_1} + \vec{V_1} \times \vec{V_2}\end{bmatrix}
    \end{flalign}
\end{center}

That's quite long isn't it? Notice that the vector part has $\vec{V_1} \times \vec{V_2}$, this means
quaternions actually are non-commutative(because vector cross products are non-commutative)! So in mathematical terms:

\begin{center}
    $Q_1Q_2 \neq Q_2Q_1$
\end{center}

But do not worry! Quaternion multiplication is actually associative:
\begin{center}
    $(Q_1Q_2)Q_3 = Q_1(Q_2Q_3)$
\end{center}

Just be careful when rearranging or calculating quaternion equations :D


\sgap
\textbf{Conjugate} of a quaternion:
\begin{center}
    $Q^* = \begin{bmatrix}w & \vec{V}\end{bmatrix}^* =\begin{bmatrix}w & -\vec{V}\end{bmatrix}$
\end{center}

The conjugate of a quaternion simply negates the vector component of a quaternion.

\sgap
\textbf{Inverse} of a quaternion:
\begin{center}
    $Q^{-1} = \frac{Q^*}{\|Q\|^2}$
\end{center}

The inverse of a quaternion applies the "inverse" rotation of a regular quaternion,
if a quaternion $Q$ is inversely applied to itself, then we would get the \textbf{identity} quaternion:
\begin{center}
    $QQ^{-1} = Q^{-1}Q = I = \begin{bmatrix}1 & \vec{0}\end{bmatrix}$
\end{center}

For a unit quaternion, due to the fact that it's magnitude being 1, it's inverse
is simply the conjugate of the quaternion!

\sgap
\textbf{Dot Product} of a quaternion:
\begin{center}
    $Q_1 \cdot Q_2 = \begin{bmatrix}w_1 & \vec{V_1}\end{bmatrix} \cdot \begin{bmatrix}w_2 & \vec{V_2}\end{bmatrix} = w_1w_2 + \vec{V_1} \cdot \vec{V_2}$
\end{center}

The quaternion dot product, similar to a vector dot product, also does element-wise multiplication, and returns a scalar
and not a quaternion.


\subsection{Quaternions as Rotations for 3D spaces(Versors)}
When we use unit quaternions for the purpose of rotating vectors, points or objects in three-dimensional spaces
, this unit quaternion is also called a "\textbf{versor}".

A versor can be defined using \textbf{Angle-Axis Form}:
\begin{center}
    $q = \begin{bmatrix}w & \vec{V}\end{bmatrix}=\begin{bmatrix}cos~\frac{\theta}{2} & sin~\frac{\theta}{2}~\hat{n}\end{bmatrix}$
\end{center}

where:
\begin{itemize}
    \item $q$: is the versor itself(or the unit quaternion).
    \item $\hat{n}$: the axis of rotation denoted as a unit vector.
    \item $\theta$: the angle of rotation around the axis denoted by $\hat{n}$.
\end{itemize}

The versor represents a pure \textbf{rotational displacement}, which can be visualized like this:
\begin{center}
    \includegraphics[width=13em, height=13em]{axis_angle_rotation}
\end{center}


You can quickly observe that:
\begin{center}
    $w = cos~\frac{\theta}{2},~ \vec{V} = sin~\frac{\theta}{2} \hat{n}$
\end{center}

You can apply a rotation to a 3D point $P=\begin{pmatrix}x & y & z\end{pmatrix}$ by first promoting this
point into a quaternion with the real component($w$) set to zero:
\begin{center}
    $P_q = \begin{bmatrix}0 & \begin{pmatrix}x & y & z\end{pmatrix}\end{bmatrix}$
\end{center}

Then apply this new promoted point quaternion in this "weird-looking" formula given a versor $Q$:
\begin{center}
    $P_q' = QP_qQ^{-1} = \begin{bmatrix}w' & \begin{pmatrix}x' & y' & z'\end{pmatrix}\end{bmatrix}$
\end{center}

The result new quaternion $P_q'$ should have the result rotated point $P'$ in it's vector component of $P_q'$.
In mathematical terms:
\begin{center}
    $P' = \begin{pmatrix}x' & y' & z'\end{pmatrix}$
\end{center}

\sgap

This formula can actually be applied multiple times consecutively to create a new rotational displacement
that is the "sum" of all of each individual rotation:
\begin{center}
    \begin{flalign}
        P_q & = b(aPa^{-1})b^{-1} \\
            & = (ba)P(a^{-1}b^{-1}) \\
            & = (ba)P(ba)^{-1}
    \end{flalign}
\end{center}

The above equation applies the versor $a$ first on $P$, then applies versor $b$ afterwards.
By rearranging the equation, it is obvious that it is equivalent to rotating the combination of two
versors $a$ and $b$.

\sgap
\subsection{Logarithm and Exponential Functions of Quaternions}

If we let $\alpha = \frac{\theta}{2}$, then we can define the logarithm function of a versor $Q$ as:
\begin{center}
    $log(Q) \equiv log(\begin{bmatrix}cos~\alpha & sin~\alpha~\hat{n}\end{bmatrix}) = \begin{bmatrix}0 & \alpha\hat{n}\end{bmatrix}$
\end{center}

As you can see, the logarithm of a versor is simply storing the half angle of rotation directly in the unit vector!
And because we know unit vectors always have a length of 1, we can obtain back the half angle $\alpha$
by simply getting the magnitude of the vector part of the quaternion(note that the magnitude
of a vector is always positive, therefore we will always lose the sign of the angle!)

But do not worry, the sign is preserved on the unit vector instead! You can see it as rotating on the opposite direction
with the reverse angle(i.e. $-\alpha$ with $-\hat{n}$).

It is also important to note the result may not be a versor / unit quaternion!

This method of storing more information in less numbers is commonly referred to as \textbf{exponential map format}.

\tgap
If a quaternion has the format $Q = \begin{bmatrix}0 & \alpha\hat{n}\end{bmatrix}$
then we can define the exponential function of a quaternion as:

\begin{center}
    $exp(Q) \equiv exp(\begin{bmatrix}0 & \alpha\hat{n}\end{bmatrix}) = \begin{bmatrix}cos~\alpha & sin~\alpha~\hat{n}\end{bmatrix}$
\end{center}

Hopefully you have noticed by now that the logarithm and exponential functions are actually inverse of each other:

\begin{center}
    $exp(log(Q)) = Q$
\end{center}

\sgap
\subsection{Exponentiation of Quaternions}
Confusingly, quaternions also have exponentiation, it is different compared to the \textbf{exponential} function. It is mathematically defined as:
\begin{center}
    $Q^t \equiv exp(t~log(Q))$
\end{center}

When applying exponentiation of $t$ to a versor that rotates $\theta$ angle around an axis, the resulting quaternion rotates
$\theta t$... With a catch! If $\theta > 180~degrees$ then what happens, is the quaternion kind of "flips to the other direction",
and rotates in the "shorter angle", here's an image that shows this:

TODO: put image here :)

So if $Q$ rotates 50 degrees, then $Q^2$ rotates 100 degrees, but $Q^4$ does not rotate 200 degrees, but rather $200 - 360 = -160~degrees$
on the other direction! (because it is shorter)

\subsection{Quaternion Difference}
TODO: Explain $\delta d = ba^{-1}$

\subsection{Quaternion Interpolation(Spherical Interpolation)}
TODO: use the above explanation of difference to derive interpolation
